\documentclass{article}
\usepackage[utf8]{inputenc}
\usepackage[T1]{fontenc}
\usepackage{geometry}
\usepackage{graphicx}
\usepackage{float}
\usepackage{hyperref}
\usepackage[english,indonesian]{babel}

\title{Optimasi Pemrosesan Template LaTeX Menggunakan Pendekatan Berbasis AI}
\author{Academic Support Agent}
\date{\today}

\begin{document}
\maketitle

\begin{abstract}

\end{abstract}

\chapter{PENDAHULUAN}


\chapter{TINJAUAN PUSTAKA}


% Original content (commented out):
% % Template LaTeX para os artigos do SBIAgro 2017 
%
% Atualizado em fevereiro de 2017 por thiago.santos@embrapa.br
%
% Versão original: SBIAgro 2015
%
\documentclass[brazil,12pt]{article}
\usepackage[utf8]{inputenc}
\usepackage{graphicx,url}
\usepackage{sbiagro2017}

\usepackage[alf]{abntcite} 

\title{Escreva aqui o título do trabalho com no máximo 15 palavras}
\author{Nome Completo Primeiro Autor\inst{1}, Nome Completo Segundo Autor\inst{2}, Nome Completo Terceiro Autor\inst{1}}

\address{Departamento de Informática, Universidade Estadual de Ponta Grossa\\
 Ponta Grossa, Paraná, Brasil\\
 primeiroautor@xxx.br, terceiroautor@zzz.br
\nextinstitute
Departamento de Agronomia, Universidade Estadual Paulista\\
Botucatu, São Paulo, Brasil\\
segundoautor@yyy.br
  }

\begin{document}

\maketitle

\thispagestyle{firstpage}

\begin{resumo} 
O resumo deve conter no máximo 250 palavras. O resumo deve ser todo
escrito em um mesmo parágrafo. As palavras-chave devem conter entre
três e cinco termos para indexação. Os termos usados não devem constar
no título. Cada palavra-chave deve iniciar com letra maiúscula e ser
seguida de vírgula. Após a última palavra-chave um ponto deve ser
usado. 
\end{resumo}
\begin{palavraschave}
Formato, Importante, Estilo.
\end{palavraschave}

\begin{abstract}
O abstract deve ser redigido em língua inglesa e precisa refletir o
conteúdo do resumo. O abstract deve ser todo escrito em um mesmo
parágrafo. As keywords devem corresponder às palavras-chave, em
inglês. Cada keyword deve iniciar com letra maiúscula e ser seguida de
vírgula. Após a última keyword um ponto deve ser usado. 
\end{abstract}
\begin{keywords}
Format, Important, Style.
\end{keywords}     

\section*{INTRODUÇÃO}
A Introdução deve ser concisa e objetiva contendo citações atuais que
apresentem relação com o assunto abordado na pesquisa. É importante
que na Introdução estejam presentes os elementos  de definição do
problema, motivação, justificativa e objetivos do trabalho. 

O primeiro parágrafo de cada seção não deve conter recuo. Nos demais
parágrafos, deve haver um recuo de 1~cm na primeira linha de cada
parágrafo. 

O texto deve estar distribuído ao longo das seções respeitando o
número total de páginas do trabalho de mínimo de 6 páginas e no máximo
10 páginas.  

A NBR 10520/2002 estabelece as condições exigidas para a apresentação
de citações em documentos técnico-científicos e acadêmicos. Nas
citações, quando o sobrenome do autor, a instituição responsável ou
título estiver incluído na sentença, este se apresenta em letras
maiúsculas/minúsculas, e quando estiverem entre parênteses, em letras
maiúsculas. 

São exemplos de citações: \citeonline{Koh95} ou \cite{Koh95}; com dois
autores ou três autores, usar \citeonline{Alm94} ou \cite{Alm94}
e \citeonline{Bra07} ou \cite{Bra07}; com mais de três autores, usar
\citeonline{CSRL2001} ou \cite{CSRL2001}. 

Em casos de existir citações indiretas de diversos documentos de
vários autores mencionados simultaneamente e que expressam a mesma
ideia, separam-se os autores por ponto e vírgula, em ordem alfabética,
independente do ano de publicação. Um exemplo desse tipo de citação é
\cite{Alm94,Koh95}.

Use o suporte \BibTeX para citações com os comandos
\texttt{citeonline} e \texttt{cite}. Exemplos de citações:
\cite{Bra07} e \citeonline{Koh95}. 

\section*{MATERIAL E MÉTODOS}
A seção de Material e Métodos  deve apresentar o material utilizado
para a condução do trabalho, bem como a metodologia empregada. Deve-se
primar pela objetividade e clareza nessa seção, para que o trabalho
desenvolvido possa ser bem compreendido e valorizado. 

As siglas devem ser usadas ao longo do trabalho da seguinte maneira:
quando aparecem pela primeira vez no texto, deve-se colocar o nome por
extenso, seguido da sigla entre parênteses. Por exemplo, podemos dizer
que esse trabalho segue as normas da Associação Brasileira de Normas
Técnicas (ABNT). 

\subsection*{Subseção 1} 
As subseções são opcionais. Quando existirem, devem ter seu título em
itálico, negrito e somente com a primeira letra em maiúsculo. O
primeiro parágrafo de cada subseção não deve conter recuo. Nos demais
parágrafos, deve haver um recuo de 1 cm na primeira linha de cada
parágrafo. 

As tabelas devem ser numeradas consecutivamente com algarismos
arábicos na parte superior. Não usar linhas verticais. As linhas
horizontais devem ser usadas para separar o título do cabeçalho e este
do conteúdo, além de uma no final da tabela. Cada dado deve ocupar uma
célula distinta. Usar espaço simples. Usar fonte 10, inclusive para o
título da tabela. Não usar negrito no cabeçalho. A
Tabela~\ref{tab:exemplo} apresenta um exemplo dessa formatação. 

\begin{table}[htb]
\caption{Titulo da tabela refletindo o conteúdo da mesma.}
\label{tab:exemplo}
\centering
\vspace{2mm}
\begin{tabular}{lccc}
\hline
Cabeçalho~1 & \multicolumn{3}{c}{Cabeçalho~2}\\\hline
            & 0--20 & 20--40 & 40--60\\\hline
Variável A & 58 & 58 & 50\\\hline
Variável B & 19 & 18 & 16\\\hline
Variável C & 16 & 16 & 16\\\hline
Variável D & 32 & 32 & 44\\\hline
Variável E (\%)$^\ddagger$ & 50 & 50 & 36\\\hline
$^\ddagger$ Variável E = 100 (C/D). 
\end{tabular}
\end{table}


Use o suporte \LaTeX para equações, conforme a Equação~\ref{eq:exemplo}:

\begin{equation}\label{eq:exemplo}
\int_0^\infty e^{-x^2} dx=\frac{\sqrt{\pi}}{2}
\end{equation}


Os gráficos, fotografias ou desenhos levarão a denominação geral de
Figura. Sua identificação, de acordo com a NBR 14724:2011, aparece na
parte superior da figura, precedida da palavra Figura, seguida de seu
número de ordem de ocorrência no texto, em algarismos arábicos,
travessão e do respectivo título. O título da figura deve ter  fonte
Times New Roman, corpo 10. Na parte inferior da figura indicar a fonte
consultada (elemento obrigatório, mesmo que seja produção do próprio
autor), legenda, notas e outras informações necessárias à sua
compreensão (se houver). A figura deve ser citada no texto (por
exemplo, Figura~\ref{fig:figura}) e inserida o mais próximo possível
do trecho a que se refere. A Figura~\ref{fig:figura} exemplifica o uso
de figuras.  

\begin{figure}[htb]
\centering
\caption{Imagem aérea de São Paulo}\label{fig:figura}
\includegraphics[width=0.5\textwidth]{figure}
\begin{center}Fonte (IBGE, 2014)\end{center}
\end{figure}

\newpage

\section*{RESULTADOS E DISCUSSÃO}
A seção de Resultados e Discussão deve apresentar os resultados do
trabalho e discuti-los em relação a outros trabalhos e resultados
associados. Assim como nas demais seções, pode conter subseções, deve
ter o primeiro parágrafo sem recuo e os demais com recuo de 1 cm na
primeira linha. 

\section*{CONCLUSÕES}
A seção de conclusões deve apresentar de forma objetiva as conclusões
obtidas com o trabalho realizado.  

\section*{AGRADECIMENTOS}
Esta seção é opcional devendo, quando existir, vir logo após as
conclusões. Nessa seção poderão ser apresentados os agradecimentos
direcionados a pessoas ou instituições, em estilo sóbrio e claro,
indicando as razões pelas quais os faz.  


\textbf{Sobre as referências:} Inicia-se com a palavra REFERÊNCIAS
(escrita em caixa alta, em negrito e centralizada).  As referências
devem ser digitadas em fonte tamanho 12, alinhadas somente à margem
esquerda do texto, em espaço simples e separadas entre si em espaço
1,5. Devem ser elaboradas conforme a ABNT 6023:2002. Deve-se dar
preferência a trabalhos publicados em periódicos científicos. Segue
exemplo a seguir. 

\bibliographystyle{abnt-alf}
\bibliography{base-referencias}

\end{document}



\end{document}